\begin{thebibliography}{00}
%第一章引用文献
\bibitem{chapterone1}
K. Samdanis and T. Taleb, ``The road beyond 5G: A vision and insight of the key technologies,'' \textit{IEEE Netw.}, vol. 34, no. 2, pp. 135-141, Feb. 2020.

\bibitem{chapterone2}
国务院, ``中华人民共和国国民经济和社会发展第十四个五年规划和2035年远景目标纲要,'' 人民出版社, 2021.

\bibitem{chapterone3}
工业和信息化部, ``十四五信息通信行业发展规划,'' 2021. [Online]. Available: https://www.miit.gov.cn/

\bibitem{chapterone4}
ITU-R, ``Framework and overall objectives of the future development of IMT for 2030 and beyond,'' Recommendation ITU-R M.2160-0, 2023.

\bibitem{chapterone5}
V. Ziegler, H. Viswanathan, H. Flinck, M. Hoffmann, V. R{\"a}is{\"a}nen, and K. H{\"a}t{\"o}nen, ``6G architecture to connect the worlds,'' \textit{IEEE Access}, vol. 8, pp. 173508-173520, Sept. 2020.

\bibitem{chapterone6}
H. Q. Ngo, A. Ashikhmin, H. Yang, E. G. Larsson, and T. L. Marzetta, ``Cell-free massive MIMO versus small cells,'' \textit{IEEE Trans. Wireless Commun.}, vol. 16, no. 3, pp. 1834--1850, Mar. 2017.

\bibitem{chapterone7}
O. T. Demir, E. Bj{\"o}rnson, and L. Sanguinetti, ``Foundations of usercentric cell-free massive MIMO,'' \textit{Found. Trends Signal Process.}, vol. 14, no. 3--4, pp. 162--472, 2021.

\bibitem{chapterone8}
G. Interdonato, E. Bj{\"o}rnson, H. Q. Ngo, P. Frenger, and E. G. Larsson, ``Ubiquitous cell-free massive MIMO communications,'' \textit{EURASIP J. Wireless Commun. Netw.}, vol. 2019, no. 1, pp. 1--13, Dec. 2019.

\bibitem{chapterone9}
S. -N. Jin, D. -W. Yue, and H. H. Nguyen, ``Spectral and energy efficiency in cell-free massive MIMO systems over correlated Rician fading,'' \textit{IEEE Syst. J.}, vol. 15, no. 2, pp. 2822--2833, Jun. 2021.

\bibitem{chapterone10}
H. Q. Ngo, L. -N. Tran, T. Q. Duong, M. Matthaiou, and E. G. Larsson, ``On the total energy efficiency of cell-free massive MIMO,'' \textit{IEEE Trans. Green Commun. Netw.}, vol. 2, no. 1, pp. 25--39, Mar. 2018.

\bibitem{chapterone11}
E. Nayebi, A. Ashikhmin, T. L. Marzetta, H. Yang, and B. D. Rao, ``Precoding and power optimization in cell-free massive MIMO systems,'' \textit{IEEE Trans. Wireless Commun.}, vol. 16, no. 7, pp. 4445--4459, Jul. 2017.

\bibitem{chapterone12}
L. D. Nguyen, T. Q. Duong, H. Q. Ngo, and K. Tourki, ``Energy efficiency in cell-free massive MIMO with zero-forcing precoding design,'' \textit{IEEE Commun. Lett.}, vol. 21, no. 8, pp. 1871--1874, Aug. 2017.

\bibitem{chapterone13}
Y. Zhang, M. Zhou, Y. Cheng, L. Yang, and H. Zhu, ``RF impairments and low-resolution ADCs for nonideal uplink cell-free massive MIMO systems,'' \textit{IEEE Syst. J.}, vol. 15, no. 2, pp. 2519-2530, Jun. 2021.

\bibitem{chapterone14}
工业和信息化部, ``信息通信行业绿色低碳发展行动计划（2022-2025年）,'' 2022. [Online]. Available: https://www.miit.gov.cn/

\bibitem{chapterone15}
J. Choi and B. L. Evans, ``Analysis of ergodic rate for transmit antenna selection in low-resolution ADC systems,'' \textit{IEEE Trans. Veh. Technol.}, vol. 68, no. 1, pp. 952-956, Jan. 2019.

\bibitem{chapterone16}
S. -N. Hong, S. Kim, and N. Lee, ``A weighted minimum distance decoding for uplink multiuser MIMO systems with low-resolution ADCs,'' \textit{IEEE Trans. Commun.}, vol. 66, no. 5, pp. 1912-1924, May 2018.

\bibitem{chapterone17}
S. Chen, J. Zhang, J. Zhang, E. Bj{\"o}rnson, and B. Ai, ``A survey on user-centric cell-free massive MIMO systems,'' 2021, \textit{arXiv:2104.13667}. [Online]. Available: http://arxiv.org/abs/2104.13667.

\bibitem{paper1}
J. Chen, L. Zhang, Y.-C. Liang, X. Kang, and R. Zhang, ``Resource allocation for wireless-powered IoT networks with short packet communication,'' \textit{IEEE Trans. Wireless Commun.}, vol. 18, no. 2, pp. 1447--1461, Feb. 2019.

\bibitem{paper2}
Y. Polyanskiy, H. V. Poor, and S. Verdu, ``Channel coding rate in the finite blocklength regime,'' \textit{IEEE Trans. Inf. Theory}, vol. 56, no. 5, pp. 2307--2359, May 2010.

\bibitem{paper3}
H. Ji, S. Park, J. Yeo, Y. Kim, J. Lee, and B. Shim, ``Ultra-reliable and low-latency communications in 5G downlink: Physical layer aspects,'' \textit{IEEE Wireless Commun.}, vol. 25, no. 3, pp. 124--130, Mar. 2018.

\bibitem{paper4}
H. Q. Ngo, A. Ashikhmin, H. Yang, E. G. Larsson, and T. L. Marzetta, ``Cell-free massive MIMO: Uniformly great service for everyone,'' in \textit{Proc. IEEE 16th Int. Workshop Signal Process. Adv. Wireless Commun. (SPAWC)}, 2015, pp. 201--205.

\bibitem{paper5}
E. Bj{\"o}rnson and L. Sanguinetti, ``Scalable cell-free massive MIMO systems,'' \textit{IEEE Trans. Commun.}, vol. 68, no. 7, pp. 4247--4261, Jul. 2020.

\bibitem{chapterone18}
J. Zhang, Y. Wei, E. Bj{\"o}rnson, Y. Han, and S. Jin, ``Performance analysis and power control of cell-free massive MIMO systems with hardware impairments,'' \textit{IEEE Access}, vol. 6, pp. 55302--55314, Sept. 2018.

\bibitem{chapterone19}
M. Xie, X. Yu, Y. Rui, K. Wang, X. Dang, and J. Zhang, ``Performance analysis for user-centric cell-free massive MIMO systems with hardware impairments and multi-antenna users,'' \textit{IEEE Trans. Wireless Commun.}, vol. 23, no. 2, pp. 1243--1259, Feb. 2024.

\bibitem{chapterone20}
B. Lim, W. J. Yun, J. Kim, and Y.-C. Ko, ``Joint pilot design and channel estimation using deep residual learning for multi-cell massive MIMO under hardware impairments,'' \textit{IEEE Trans. Veh. Technol.}, vol. 71, no. 7, pp. 7599--7612, Jul. 2022.

\bibitem{chapterone21}
Y. Zhang, Q. Zhang, H. Hu, L. Yang, and H. Zhu, ``Cell-free massive MIMO systems with non-ideal hardware: Phase drifts and distortion noise,'' \textit{IEEE Trans. Veh. Technol.}, vol. 70, no. 11, pp. 11604-11618, Nov. 2021.

\bibitem{chapterone22}
S. Jacobsson, U. Gustavsson, G. Durisi, and C. Studer, ``Massive MU-MIMO-OFDM uplink with hardware impairments: Modeling and analysis,'' in \textit{Proc. Asilomar Conf. Signals, Syst. Comput. (ACSSC)}, Pacific Grove, CA, USA, 2018, pp. 1829-1835.

\bibitem{chapterone23}
E. Bj{\"o}rnson, J. Hoydis, M. Kountouris, and M. Debbah, ``Massive MIMO systems with non-ideal hardware: Energy efficiency, estimation, and capacity limits,'' \textit{IEEE Trans. Inf. Theory.}, vol. 60, no. 11, pp. 7112-7139, Nov. 2014.

\bibitem{chapterone24}
E. Bj{\"o}rnson, J. Hoydis, and L. Sanguinetti, ``Massive MIMO networks: Spectral, energy, and hardware efficiency,'' \textit{Found. Trends Signal Process.}, 2017.

\bibitem{chapterone25}
Q. Sun, X. Ji, Z. Wang, X. Chen, Y. Yang, J. Zhang, and K. Wong, ``Uplink performance of hardware-impaired cell-free massive MIMO with multi-antenna users and superimposed pilots,'' \textit{IEEE Trans. Commun.}, vol. 71, no. 11, pp. 6711-6726, Nov. 2023.

\bibitem{chapterone26}
X. Hu, C. Zhong, X. Chen, W. Xu, H. Lin, and Z. Zhang, ``Cell-free massive MIMO systems with low resolution ADCs,'' \textit{IEEE Trans. Commun.}, vol. 67, no. 10, pp. 6844-6857, Oct. 2019.

\bibitem{chapterone27}
Y. Zhang, M. Zhou, X. Qiao, H. Cao, and L. Yang, ``On the performance of cell-free massive MIMO with low-resolution ADCs,'' \textit{IEEE Access}, vol. 7, pp. 117968-117977, Aug. 2019.

\bibitem{chapterone28}
M. Zhou, L. Yang, and H. Zhu, ``Sum-SE for multigroup multicast cell-free massive MIMO with multi-antenna users and low-resolution DACs,'' \textit{IEEE Wireless Commun. Lett.}, vol. 10, no. 8, pp. 1702-1706, Aug. 2021.

\bibitem{chapterone29}
M. Zhou, J. Li, M. Xie, C. Zhang, and J. Yuan, ``Average sum rate of D2D underlaid multigroup multicast cell-free massive MIMO with multi-antenna users,'' \textit{IEEE Wireless Commun. Lett.}, vol. 12, no. 1, pp. 60-64, Jan. 2023.

\bibitem{chapterone30}
F. Meng, P. Chen, L. Wu, and J. Cheng, ``Power allocation in multi-user cellular networks: Deep reinforcement learning approaches,'' \textit{IEEE Trans. Wireless Commun.}, vol. 19, no. 10, pp. 6255-6267, Oct. 2020.

\bibitem{chapterone31}
R. S. Sutton and A. G. Barto, \textit{Reinforcement learning: An introduction.} Cambridge, MA, USA: MIT Press, 2018.

\bibitem{chapterone32}
Y. Zhao, F. Zhang, Z. Zhou, and G. Hu, ``A dynamic power allocation approach for downlink cell-free massive MIMO with graph neural network,'' \emph{IEEE Trans. Veh. Technol.}, doi: 10.1109/TVT.2024.

\bibitem{chapterone33}
Z. Liu, J. Zhang, E. Shi, Z. Liu, D. Niyato, B. Ai, and X. S. Shen, ``Graph neural network meets multi-agent reinforcement learning: Fundamentals, applications, and future directions,'' \emph{IEEE Wireless Commun.}, vol. 31, no. 6, pp. 39-47, Dec. 2024.

\bibitem{chapterone34}
Y. Shen, J. Zhang, S. H. Song, and K. B. Letaief, ``Graph neural networks for wireless communications: From theory to practice,'' \emph{IEEE Trans. Wireless Commun.}, vol. 22, no. 5, pp. 3554--3569, May 2023.

\bibitem{chapterone35}
Y. Gu, C. She, S. Bi, Z. Quan, and B. Vucetic, ``Graph neural network for distributed beamforming and power control in massive URLLC networks,'' \textit{IEEE Trans. Wireless Commun.}, vol. 23, no. 8, pp. 9099--9112, Aug. 2024.

\bibitem{chapterone36}
L. U. Khan et al., ``Federated learning for edge networks: Resource optimization and incentive mechanism,'' \textit{IEEE Commun. Mag.}, vol. 58, no. 10, pp. 88--93, Oct. 2020.

\bibitem{chapterone37}
Y. LeCun, Y. Bengio, and G. Hinton, ``Deep learning,'' \textit{Nature}, vol. 521, no. 7553, pp. 436--444, 2015.

\bibitem{chapterone38}
H. Sun, X. Chen, Q. Shi, M. Hong, X. Fu, and N. D. Sidiropoulos, ``Learning to optimize: Training deep neural networks for interference management,'' \textit{IEEE Trans. Signal. Process.}, vol. 66, no. 20, pp. 5438--5453, Oct. 2018.

\bibitem{chapterone39}
R. Sun, N. Cheng, C. Li, F. Chen, and W. Chen, ``Knowledge-driven deep learning paradigms for wireless network optimization in 6G,'' \textit{IEEE Netw.}, vol. 38, no. 2, pp. 70--78, Mar. 2024.

\bibitem{chapterone40}
F. Liang, C. Shen, W. Yu, and F. Wu, ``Towards optimal power control via ensembling deep neural networks,'' \textit{IEEE Trans. Commun.}, vol. 68, no. 3, pp. 1760--1776, Mar. 2020.

\bibitem{chapterone41}
W. Lee and R. Schober, ``Deep learning-based resource allocation for device-to-device communication,'' \textit{IEEE Trans. Wireless Commun.}, vol. 21, no. 7, pp. 5235--5250, Jul. 2022.

\bibitem{chapterone42}
A. R. Barron, ``Universal approximation bounds for superpositions of a sigmoidal function,'' \textit{IEEE Trans. Inf. Theory.}, vol. 39, no. 3, pp. 930--945, May 1993.

\bibitem{chapterone43}
R. Dong, C. She, W. Hardjawana, Y. Li, and B. Vucetic, ``Deep learning for radio resource allocation with diverse quality-of-service requirements in 5G,'' \textit{IEEE Trans. Wireless Commun.}, vol. 20, no. 4, pp. 2309--2324, Apr. 2021.

%第二章引用文献
\bibitem{chaptertwo1}
Y. Polyanskiy, H. V. Poor, and S. Verd\'{u}, ``Channel coding rate in the finite blocklength regime,'' \emph{IEEE Trans. Inf. Theory}, vol. 56, no. 5, pp. 2307--2359, May 2010.

\bibitem{chaptertwo2}
C. Sun, C. She, C. Yang, and T. Q. S. Quek, ``Optimizing resource allocation in the short blocklength regime for ultra-reliable and low-latency communications,'' \emph{IEEE Trans. Wireless Commun.}, vol. 18, no. 1, pp. 402--415, Jan. 2019.

\bibitem{chaptertwo3}
C. She, C. Yang, and T. Q. S. Quek, ``Cross-layer optimization for ultra-reliable and low-latency radio access networks,'' \emph{IEEE Trans. Wireless Commun.}, vol. 17, no. 1, pp. 127--141, Jan. 2018.

\bibitem{chaptertwo4}
3GPP TS 38.913, ``Study on scenarios and requirements for next generation access technologies,'' v15.0.0, 2018.

\bibitem{chaptertwo5}
M. Bennis, M. Debbah, and H. V. Poor, ``Ultrareliable and low-latency wireless communication: Tail, risk, and scale,'' \emph{Proc. IEEE}, vol. 106, no. 10, pp. 1834--1853, Oct. 2018.

\bibitem{chaptertwo6}
T. N. Kipf and M. Welling, ``Semi-supervised classification with graph convolutional networks,'' in \emph{Proc. ICLR}, Toulon, France, Apr. 2017.

\bibitem{chaptertwo7}
D. I. Shuman, S. K. Narang, P. Frossard, A. Ortega, and P. Vandergheynst, ``The emerging field of signal processing on graphs: Extending high-dimensional data analysis to networks and other irregular domains,'' \emph{IEEE Signal Process. Mag.}, vol. 30, no. 3, pp. 83--98, May 2013.

\bibitem{chaptertwo8}
J. Bruna, W. Zaremba, A. Szlam, and Y. LeCun, ``Spectral networks and locally connected networks on graphs,'' in \emph{Proc. ICLR}, Banff, Canada, Apr. 2014.

\bibitem{chaptertwo9}
M. Defferrard, X. Bresson, and P. Vandergheynst, ``Convolutional neural networks on graphs with fast localized spectral filtering,'' in \emph{Proc. NIPS}, Barcelona, Spain, Dec. 2016, pp. 3844--3852.

\bibitem{chaptertwo10}
P. Veli\v{c}kovi\'{c}, G. Cucurull, A. Casanova, A. Romero, P. Li\`{o}, and Y. Bengio, ``Graph attention networks,'' in \emph{Proc. ICLR}, Vancouver, Canada, May 2018.

\bibitem{chaptertwo11}
R. S. Sutton and A. G. Barto, \emph{Reinforcement Learning: An Introduction}, 2nd ed. Cambridge, MA, USA: MIT Press, 2018.

\bibitem{chaptertwo12}
V. Mnih et al., ``Human-level control through deep reinforcement learning,'' \emph{Nature}, vol. 518, no. 7540, pp. 529--533, Feb. 2015.

\bibitem{chaptertwo13}
V. Mnih et al., ``Playing Atari with deep reinforcement learning,'' in \emph{Proc. NIPS Deep Learning Workshop}, Lake Tahoe, USA, Dec. 2013.

\bibitem{chaptertwo14}
J. Schulman, F. Wolski, P. Dhariwal, A. Radford, and O. Klimov, ``Proximal policy optimization algorithms,'' arXiv preprint arXiv:1707.06347, 2017.

\bibitem{chaptertwo15}
T. P. Lillicrap et al., ``Continuous control with deep reinforcement learning,'' in \emph{Proc. ICLR}, San Juan, Puerto Rico, May 2016.

\bibitem{chaptertwo16}
D. Silver et al., ``Deterministic policy gradient algorithms,'' in \emph{Proc. ICML}, Beijing, China, Jun. 2014, pp. 387--395.

%第三章引用文献
\bibitem{chapterthree1}
U. Gustavsson, C. Sanch\'{e}z-Perez, T. Eriksson, F. Athley, G. Durisi, and P. Landin, ``On the impact of hardware impairments on massive MIMO,'' in \emph{Proc. IEEE Globecom Workshops (GC Wkshps)}, Austin, TX, USA, 2014, pp. 294--300.

\bibitem{chapterthree2}
A. Mohammadi and F. M. Ghannouchi, \textit{RF transceiver design for MIMO wireless communications.} Springer Science \& Business Media, 2012, vol. 145.

\bibitem{chapterthree3}
D. R{\"o}nnow and P. H{\"a}ndel, ``Nonlinear distortion noise and linear attenuation in MIMO systems---Theory and application to multiband transmitters,'' \emph{IEEE Trans. Signal Process.}, vol. 67, no. 20, pp. 5203--5212, Oct. 2019.

\bibitem{chapterthree4}
P. H{\"a}ndel and D. R{\"o}nnow, ``Dirty MIMO transmitters: Does it matter?'' \emph{IEEE Trans. Wireless Commun.}, vol. 17, no. 8, pp. 5425--5436, Aug. 2018.

\bibitem{chapterthree5}
{\"O}. T. Demir and E. Bj{\"o}rnson, ``Channel estimation in massive MIMO under hardware non-linearities: Bayesian methods versus deep learning,'' \emph{IEEE Open J. Commun. Soc.}, vol. 1, pp. 109-124, Dec. 2020.

\bibitem{chapterthree6}
M. Jadidi, A. M. Khoueini, A. Mohammadi, and V. Meghdadi, ``Performance analysis and power allocation for uplink cell-free massive MIMO system with nonlinear power amplifier,'' \emph{IEEE Trans. Commun.}, vol. 72, no. 9, pp. 5473-5485, Sept. 2024.

\bibitem{chapterthree7}
J. Zhang, L. Dai, S. Sun, and Z. Wang, ``On the spectral efficiency of massive MIMO systems with low-resolution ADCs,'' \textit{IEEE Commun. Lett.}, vol. 20, no. 5, pp. 842-845, May 2016.

\bibitem{chapterthree8}
J. Zhang, L. Dai, Z. He, S. Jin, and X. Li, ``Performance analysis of mixed-ADC massive MIMO systems over rician fading channels,'' \textit{IEEE J. Sel. Areas Commun.}, vol. 35, no. 6, pp. 1327-1338, Jun. 2017.

\bibitem{chapterthree9}
J. Yuan, S. Jin, C. -K. Wen, and K. -K. Wong, ``The distributed MIMO scenario: Can ideal ADCs be replaced by low-resolution ADCs?,'' \textit{IEEE Wireless Commun. Lett.}, vol. 6, no. 4, pp. 470-473, Aug. 2017.

\bibitem{chapterthree10}
H. Moazzen, A. Mohammadi, and M. Majidi, ``Performance analysis of linear precoded MU-MIMO-OFDM systems with nonlinear power amplifiers and correlated channel,'' \textit{IEEE Trans. Commun.}, vol. 67, no. 10, pp. 6753-6765, Oct. 2019.

%第四章引用文献（无新引用，使用前面章节的引用）

\end{thebibliography}
